
%%%%%%%%%%%%%%%%%%%%%%%%%%%%%%%%%%%%%%%%%%%%%%%%%%%%%%%%%%%%%%%%%%%%%%%%%%%%%%%%
%%%% DOCUMENT SETUP                                                         %%%%
%%%%%%%%%%%%%%%%%%%%%%%%%%%%%%%%%%%%%%%%%%%%%%%%%%%%%%%%%%%%%%%%%%%%%%%%%%%%%%%%

% Sub-enumerates include the enumerate.
%\usepackage{enumitem}
%\setlist[enumerate,2]{label*=\arabic*.}
\usepackage{multicol}

%\pretocmd{\tableofcontents}{\begin{minipage}{\textwidth}}{}{}
%\apptocmd{\tableofcontents}{\end{minipage}}{}{}


%%%%%%%%%%%%%%%%%%%%%%%%%%%%%%%%%%%%%%%%%%%%%%%%%%%%%%%%%%%%%%%%%%%%%%%%%%%%%%%%
%%%% FONTS                                                                  %%%%
%%%%%%%%%%%%%%%%%%%%%%%%%%%%%%%%%%%%%%%%%%%%%%%%%%%%%%%%%%%%%%%%%%%%%%%%%%%%%%%%
\usefonttheme{professionalfonts}

\usepackage[cal=txupr]{mathalpha}
\usepackage{newtxtext, dsfont}
\usepackage[nopatch=footnote]{microtype}
\usepackage[misc]{ifsym}
\usepackage[normalem]{ulem}

% Source - https://tex.stackexchange.com/a/146531
% Posted by pmk, modified by community. See post 'Timeline' for change history
% Retrieved 2026-08-25, License - CC BY-SA 4.0
\renewcommand{\ULthickness}{1.5pt}
\renewcommand<>{\sout}[1]{\alt#2{\beameroriginal{\sout}{#1}}{#1}}




%%%%%%%%%%%%%%%%%%%%%%%%%%%%%%%%%%%%%%%%%%%%%%%%%%%%%%%%%%%%%%%%%%%%%%%%%%%%%%%%
%%%% VISUALS                                                                %%%%
%%%%%%%%%%%%%%%%%%%%%%%%%%%%%%%%%%%%%%%%%%%%%%%%%%%%%%%%%%%%%%%%%%%%%%%%%%%%%%%%
\usepackage{caption, tikz, tikz-3dplot, xcolor, wrapfig, nicematrix, circuitikz}
\usetikzlibrary{positioning, graphs, calc, decorations, decorations.markings, patterns, patterns.meta, overlay-beamer-styles, through}


% Captions.
\captionsetup{format=plain, font={it, footnotesize}}
\setlength{\belowcaptionskip}{-1em}



%%%%%%%%%%%%%%%%%%%%%%%%%%%%%%%%%%%%%%%%%%%%%%%%%%%%%%%%%%%%%%%%%%%%%%%%%%%%%%%%
%%%% COLORS                                                                 %%%%
%%%%%%%%%%%%%%%%%%%%%%%%%%%%%%%%%%%%%%%%%%%%%%%%%%%%%%%%%%%%%%%%%%%%%%%%%%%%%%%%

\definecolor{almond}{rgb}{0.94, 0.87, 0.8}
\definecolor{antiquewhite}{rgb}{0.98, 0.92, 0.84}

% IBM Design Library palette. (Colorblind-safe) %% https://tinyurl.com/27azp7vh
\definecolor{ibm-gold}{HTML}{ffb000}
\definecolor{ibm-orange}{HTML}{fe6100}
\definecolor{ibm-pink}{HTML}{dc267f}
\definecolor{ibm-violet}{HTML}{785ef0}
\definecolor{ibm-blue}{HTML}{648fff}


% Paul Tol's high-contrast palette. Colors remain distinct in monochrome. %% https://tinyurl.com/2b9mab2z
\definecolor{tol-hc-yellow}{RGB}{221,170,51}
\definecolor{tol-hc-red}{RGB}{187,85,102}
\definecolor{tol-hc-blue}{RGB}{0,68,136}


% Paul Tol's medium-contrast palette. Colors remain distinct in monochrome. (Colorblind-safe) %% https://tinyurl.com/2b9mab2z
\definecolor{tol-mc-lightyellow}{RGB}{238,204,102}
\definecolor{tol-mc-lightred}{RGB}{238,153,170}
\definecolor{tol-mc-lightblue}{RGB}{102,153,204}
\definecolor{tol-mc-darkyellow}{RGB}{153,119,0}
\definecolor{tol-mc-darkred}{RGB}{153,68,85}
\definecolor{tol-mc-darkblue}{RGB}{0,68,136}


% Paul Tol's vibrant palette. (Colorblind-safe) %% https://tinyurl.com/2b9mab2z
\definecolor{tol-vibrant-blue}{RGB}{0,119,187}
\definecolor{tol-vibrant-cyan}{RGB}{51,187,238}
\definecolor{tol-vibrant-teal}{RGB}{0,153,136}
\definecolor{tol-vibrant-orange}{RGB}{238,119,51}
\definecolor{tol-vibrant-red}{RGB}{204,51,17}
\definecolor{tol-vibrant-magenta}{RGB}{238,51,119}
\definecolor{tol-vibrant-grey}{RGB}{187,187,187}


% Okabe and Ito's palette. (Colorblind-safe) %% https://tinyurl.com/22lwpexl
\definecolor{oi-orange}{RGB}{230,159,0}
\definecolor{oi-cyan}{RGB}{86,180,233}
\definecolor{oi-green}{RGB}{0,158,115}
\definecolor{oi-yellow}{RGB}{240,228,66}
\definecolor{oi-blue}{RGB}{0,114,178}
\definecolor{oi-red}{RGB}{213,94,0}
\definecolor{oi-pink}{RGB}{204,121,167}

% ArXiV color (RIP).
\definecolor{arxiv}{HTML}{AE2A24}


% Define a color theme; use Paul Tol's dark, high-contrast palette.
\def\primarycolor{tol-hc-red}
\def\secondarycolor{tol-hc-blue}
\def\tertiarycolor{tol-hc-yellow}


%%%%%%%%%%%%%%%%%%%%%%%%%%%%%%%%%%%%%%%%%%%%%%%%%%%%%%%%%%%%%%%%%%%%%%%%%%%%%%%%
%%%% ENVIRONMENTS                                                           %%%%
%%%%%%%%%%%%%%%%%%%%%%%%%%%%%%%%%%%%%%%%%%%%%%%%%%%%%%%%%%%%%%%%%%%%%%%%%%%%%%%%
\usepackage{amsthm}
\setbeamertemplate{theorems}[numbered]
\usepackage[breakable,skins]{tcolorbox}
\usepackage[thmtools-compat]{keytheorems}

% Define standard heights, widths, etc.
\def\spaceabove{0.5\parskip}		% thmtools above
\def\spacebelow{0.25\parskip}		% thmtools below
\def\boxpadding{0.5em}

\def\postheadhook{\strut}
\def\prefoothook{\looseness=-1\strut}

\def\rulewidth{1pt}					% rule width
\def\transparency{15}				% transparency

\setbeamertemplate{itemize item}{\tiny\color{black}$\blacktriangleright$}


% Default tcolorbox styles that do some of the inheritance for us.
\tcbset{
	default/.style={
		enhanced,
		frame hidden,
		sharp corners,
		left=0.5em,
		right=0.5em,
		bottom=0.25em,
		top=0.25em,
%		before skip=\spaceabove,
%		after skip=\spacebelow,
		left skip=0pt,
		right skip=0pt,
		boxsep=0pt,
		width=0.75\textwidth,
		boxrule=0pt
	},
	break/.style={
		default,
		breakable
	},
	nobreak/.style={
		default
	},
	wide/.style={
		width=0.9\textwidth
	}
}

% Default keytheorems environment that also does some inheritance
\declarekeytheoremstyle{default}{
	postheadhook=\postheadhook,
	prefoothook=\prefoothook,
	bodyfont=\normalfont,
%	margin=0.25em
}

\declarekeytheoremstyle{compact}{
	inherit-style=default,
	spaceabove=\spaceabove,
	spacebelow=\spacebelow,
	bodyfont=\itshape
}


\declaretheoremstyle[
	inherit-style=default,
	tcolorbox-no-titlebar={
		break,
		colback=gray!\transparency
	}
]{shaded}

\declaretheoremstyle[
	inherit-style=default,
	tcolorbox-no-titlebar={
		break,
		wide,
		colback=gray!\transparency
	}
]{shaded-wide}

\declaretheoremstyle[
	inherit-style=default,
	bodyfont=\itshape,
	tcolorbox-no-titlebar={
		break,
		colback=gray!\transparency
	}
]{shaded-it}


\declaretheoremstyle[
	inherit-style=default,
	tcolorbox-no-titlebar={
		nobreak,
		colback=gray!\transparency
	}
]{shaded-nobreak}

\declaretheoremstyle[
	inherit-style=default,
	tcolorbox-no-titlebar={
		nobreak,
		wide,
		colback=gray!\transparency
	}
]{shaded-nobreak-wide}



\declaretheoremstyle[
	inherit-style=default,
	bodyfont=\itshape,
	tcolorbox-no-titlebar={
		nobreak,
		colback=gray!\transparency
	}
]{shaded-nobreak-it}


\declaretheoremstyle[
	inherit-style=default,
	tcolorbox-no-titlebar={
		break,
		colback=almond!\transparency
	}
]{shaded-almond}

\declaretheoremstyle[
	inherit-style=default,
	tcolorbox-no-titlebar={
		break,
		wide,
		colback=almond!\transparency
	}
]{shaded-almond-wide}


\declaretheoremstyle[
	inherit-style=default,
	tcolorbox-no-titlebar={
		nobreak,
		colback=almond!\transparency
	}
]{shaded-almond-nobreak}


\declaretheoremstyle[
	inherit-style=default,
	tcolorbox-no-titlebar={
		nobreak,
		colback=\primarycolor!\transparency
	}
]{shaded-primary-nobreak}

\declaretheoremstyle[
	inherit-style=default,
	tcolorbox-no-titlebar={
		nobreak,
		colback=\secondarycolor!\transparency
	}
]{shaded-secondary-nobreak}

\declaretheoremstyle[
	inherit-style=default,
	tcolorbox-no-titlebar={
		nobreak,
		colback=\tertiarycolor!\transparency
	}
]{shaded-tertiary-nobreak}


\declaretheoremstyle[
	inherit-style=default,
	tcolorbox-no-titlebar={
		nobreak,
		wide,
		colback=\tertiarycolor!\transparency
	}
]{shaded-tertiary-nobreak-wide}


\declaretheorem[%
	numberwithin=subsection,
	style=compact
]{theorem}

\declaretheorem[%
	numbered=no,
	style=shaded-primary-nobreak,
]{lemma, definition, corollary, proposition, conjecture}

\declaretheorem[%
	numbered=no,
	style=shaded-nobreak,
	title=Theorem
]{bigtheorem}

\declaretheorem[%
	numbered=no,
	style=shaded-nobreak,
	title=Corollary
]{bigcorollary}

\declaretheorem[%
	numbered=no,
	style=shaded-secondary-nobreak,
	title=Definition
]{bigdefinition}

\declaretheorem[%
	sharenumber=theorem,
	style=shaded-nobreak
]{example}

\declaretheorem[%
	numbered=no,
	style={shaded-wide}
]{question}

\declaretheorem[%
%	sharenumber=question,
	numbered=no,
	style={shaded-almond-wide}
]{goal}

\declaretheorem[%
%	sharenumber=question,
	numbered=no,
	style={shaded-almond},
	title={Goal}
]{narrowgoal}

\declaretheorem[%
	numbered=no,
	style=shaded-primary-nobreak,
	title={Author's Note}
]{authorsnote}

\declaretheorem[%
	numbered=no,
	style=shaded-tertiary-nobreak,
	title={Note}
]{readersnote}

\declaretheorem[%
	numbered=no,
	style=shaded-tertiary-nobreak-wide,
	title={Note}
]{widenote}


%
% Algorithm environment.
\tcbuselibrary{breakable,skins}
\newtcolorbox[auto counter]{boxedalgorithm}[3]{%
	label=#2,
	%
	enhanced,
	sharp corners,
	frame hidden,
	%
	center,
	width=0.75\linewidth,
	top=0.5em,
	bottom=0.5em,
	left=0.5em,
	right=0.5em,
	%
	before skip=1em,
	after skip=1.5em,
	%
	attach title to upper={#3},
	coltitle=black,
%	title={\uline{\textbf{Algorithm \thetcbcounter}{#1}\hskip0pt plus 1filll}},
	title={\underline{\makebox[\linewidth][l]{\textbf{Algorithm \thetcbcounter}{#1}}}},
	%
	breakable
}
%
%\newtcolorbox[auto counter]{headlessboxedalgorithm}{%
%	%
%	enhanced,
%	sharp corners,
%	frame hidden,
%	%
%	center,
%	width=\linewidth,
%	top=0.5em,
%	bottom=0.5em,
%	left=0.5em,
%	right=0.5em,
%	%
%	before skip=0.5em,
%	after skip=0.5em,
%	%
%	attach title to upper,
%	coltitle=black,
%	%
%	breakable
%}
%
%
%% List-like algorithm setup.
%\newlist{enumalgorithm}{enumerate}{2}
%\setlist[enumalgorithm,1]{nosep, topsep=0.25em, itemsep=2.5pt, label=(\arabic*), leftmargin=*}
%\setlist[enumalgorithm,2]{topsep=2.5pt, parsep=0pt, itemsep=2.5pt, label=(\alph*), ref=(\arabic{enumalgorithmi}\alph*)}
%
%
% Alternate description environment.
%\newlist{descriptor}{description}{1}
%\setlist[descriptor]{labelindent=0.25in,labelwidth=0.25in,leftmargin=!,align=left,itemsep=0pt,topsep=0.5em,rightmargin=\leftmargin,font={\bf\textit}}
%
%\newlist{widedescriptor}{description}{1}
%\setlist[widedescriptor]{topsep=0pt,labelindent=0.25in,labelwidth=0.25in,leftmargin=!,align=left,itemsep=0pt,topsep=0.5em,font={\bf\textit}}
%
%
% Inset algorithm environment that wraps around boxedalgorithm and enumalgorithm.
\makeatletter

\define@key{insetalgorithm@keys}{name}{\def\insetalgorithm@name{#1}}%
\define@key{insetalgorithm@keys}{label}{\def\insetalgorithm@label{#1}}%
\define@key{insetalgorithm@keys}{description}{\def\insetalgorithm@description{#1}}%

\newenvironment{insetalgorithm}[1][]
	{%
		\settowidth{\leftmargini}{\usebeamertemplate{itemize item}}
		\addtolength{\leftmargini}{\labelsep}
		\setkeys{insetalgorithm@keys}{name={},label=algorithm,description={},#1}
		\begin{boxedalgorithm}{
			\ifthenelse{\equal{\insetalgorithm@name}{}}{}{%
				\textit{(\insetalgorithm@name)}%
			}.%
		}{\insetalgorithm@label}{%
			\ifthenelse{\equal{\insetalgorithm@description}{}}{\\[-1em]%
			}{%
                \\[0.5em]\insetalgorithm@description%
			}
		}
			\begin{enumerate}%
	}
	{%
			\end{enumerate}
		\end{boxedalgorithm}%
	}

\makeatother


% Define a few universal graphical constants.

% Variables
\def\width{4}
\def\height{\width}
\def\depth{\width}
\def\scale{0.85}
\def\edgewidth{2pt}

% TikZ styles
\tikzset{
	vertex/.style={
		draw,
		line width={\edgewidth},
		black,
		circle,
		inner sep=2pt,
		fill=white
	},
	%
	edge/.style={
		line width={\edgewidth},
		line cap=round,
		black
	},
	thickie/.style={
		line cap=round,
		line width={1.5*\edgewidth},
	},
	root/.style={
		thickie,
		\primarycolor
	},
	%
	square/.style={
%		draw=none,
		fill=\secondarycolor!\transparency
	},
	%
	forward/.style={
		solid,
		currarrow,
		pos=0.5,
		sloped,
		scale=0.5
	},
	%
	cycle/.style={
		\tertiarycolor,
		line width=1.5*\edgewidth
	},
	%
	boundary/.style={
		\primarycolor,
		line width=1.5*\edgewidth
	},
	forest/.style={line width=\edgewidth},
	missing/.style={line width=\edgewidth, dashed, gray},
	rootvertex/.style={vertex, fill=\primarycolor},
	%
	% squares and stuff
	dualedge/.style={thickie,\primarycolor},
	dualsquare/.style={edge},
	pics/side/.style args={#1}{
		code={
			\draw[dualsquare,#1] (0,0,0) -- (0,1,0) -- (0,1,1) -- (0,0,1) -- cycle;
		}
	},
	pics/bottom/.style args={#1}{
		code={
			\draw[dualsquare,#1] (0,0,0) -- (1,0,0) -- (1,0,1) -- (0,0,1) -- cycle;
		}
	},
	pics/back/.style args={#1}{
		code={
			\draw[dualsquare,#1] (0,0,0) -- (0,1,0) -- (1,1,0) -- (1,0,0) -- cycle;
		}
	},
	pics/bottomt/.style args={#1}{
		code={
			\draw[dualsquare] (0,0,0) -- (1,0,0);
			\draw[dualedge] (#1,-{#1},#1) -- (#1,#1,#1);
			\draw[dualsquare] (0,0,0) -- (0,0,1) -- (1,0,1) -- (1,0,0);
		}
	},
	pics/backt/.style args={#1}{
		code={
			\draw[dualedge] (#1,#1,-{#1}) -- (#1,#1,#1);
			\draw[dualsquare] (0,0,0) -- (0,1,0) -- (1,1,0) -- (1,0,0) -- cycle;
		}
	},
	pics/sidet/.style args={#1}{
		code={
			\draw[dualsquare] (0,0,0) -- (0,1,0);
			\draw[dualedge] (-{#1},#1,#1) -- (#1,#1,#1);
			\draw[dualsquare] (0,0,0) -- (0,0,1) -- (0,1,1) -- (0,1,0);
		}
	}
}



%%%%%%%%%%%%%%%%%%%%%%%%%%%%%%%%%%%%%%%%%%%%%%%%%%%%%%%%%%%%%%%%%%%%%%%%%%%%%%%%
%%%% BIBLIOGRAPHY                                                           %%%%
%%%%%%%%%%%%%%%%%%%%%%%%%%%%%%%%%%%%%%%%%%%%%%%%%%%%%%%%%%%%%%%%%%%%%%%%%%%%%%%%
\usepackage[
	backend=biber,
	style=numeric-comp,
	sorting=noneyear,
	sortcites,
	maxnames=3,
	date=year
]{biblatex}

% Italicized titles
\DeclareFieldFormat[article,inproceedings,book,online,incollection]{title}{\emph{#1}}

% Remove extraneous fields.
\AtEveryBibitem{\clearfield{pages}}
\AtEveryBibitem{\clearfield{volume}}
\AtEveryBibitem{\clearfield{publisher}}
\AtEveryBibitem{\clearfield{editor}}
\AtEveryBibitem{\clearfield{series}}
\AtEveryBibitem{\clearfield{number}}
\AtEveryBibitem{\clearfield{isbn}}

\AtEveryBibitem{\clearlist{series}}
\AtEveryBibitem{\clearlist{publisher}}
\AtEveryBibitem{\clearlist{editor}}

% Re-format useful fields.
\DeclareFieldFormat*{url}{}
\DeclareFieldFormat[software]{url}{\mkbibacro{URL}\addcolon\space\url{#1}}

\DeclareFieldFormat*{doi}{}
\DeclareFieldFormat[online,software]{doi}{\mkbibacro{DOI}\addcolon\space\url{#1}}


% Smaller (list) label sep, itemsep, font size, multiple citations.
\setlength{\biblabelsep}{\labelsep}
\renewcommand*{\bibitemsep}{0pt}
\renewcommand*{\bibfont}{\footnotesize}
\renewcommand*{\multicitedelim}{\addcomma\nobreakspace}

% Sort references by year within the citation.
\DeclareSortingTemplate{noneyear}{
	\sort{\citeorder}
	\sort{\field{year}}
}

\setbeamertemplate{bibliography entry article}{}
\setbeamertemplate{bibliography entry title}{}
\setbeamertemplate{bibliography entry location}{}
\setbeamertemplate{bibliography entry note}{}

\setbeamercolor*{bibliography item}{fg=black}
\setbeamercolor*{bibliography entry title}{fg=black}
\setbeamercolor*{bibliography entry author}{fg=black}
\setbeamercolor*{bibliography entry location}{fg=black}
\setbeamercolor*{bibliography entry note}{fg=black}

\addbibresource{main.bib}



%%%%%%%%%%%%%%%%%%%%%%%%%%%%%%%%%%%%%%%%%%%%%%%%%%%%%%%%%%%%%%%%%%%%%%%%%%%%%%%%
%%%% SPECIAL MATHEMATICS TYPESETTING                                        %%%%
%%%%%%%%%%%%%%%%%%%%%%%%%%%%%%%%%%%%%%%%%%%%%%%%%%%%%%%%%%%%%%%%%%%%%%%%%%%%%%%%
\usepackage{amsmath, amssymb, nicefrac, mathtools, bm}

\renewcommand{\qed}{\hfill$\blacksquare$}

% Allow aligned equations to break across pages.
\allowdisplaybreaks

\let\epsilon\varepsilon

\providecommand{\N}{\mathds{N}}					% naturals
\providecommand{\F}{\mathds{F}}					% arbitrary field
\providecommand{\R}{\mathds{R}}					% reals
\providecommand{\C}{\mathds{C}}					% complex
\providecommand{\T}{\mathds{T}}					% torus?
\providecommand{\Z}{\mathds{Z}}					% integers
\providecommand{\Q}{\mathds{Q}}					% rationals
\providecommand{\RP}{\mathds{R}\mathrm{P}}		% projective plane

\let\P\undefined
\let\H\undefined

\DeclareMathOperator{\P}{\mathbf{P}}				% probability measure
\DeclareMathOperator{\E}{\mathbf{E}}				% expectation
\DeclareMathOperator{\1}{\mathbf{1}}				% indicator
\DeclareMathOperator{\Uniform}{\mathbf{Uniform}}	% uniform distribution

\DeclareMathOperator{\Laplacian}{\mathcal L}	% Laplacian

\let\im\undefined
\let\ker\undefined
\DeclareMathOperator{\ker}{\mathrm{ker}}
\DeclareMathOperator{\im}{\mathrm{im}}			% algebraic image
\DeclareMathOperator{\Hom}{\mathrm{Hom}}		% Hom functor
\DeclareMathOperator{\Ext}{\mathrm{Ext}}		% Ext functor
\DeclareMathOperator{\Tor}{\mathrm{Tor}}		% Tor functor
\DeclareMathOperator{\rank}{\mathrm{rank}}		% algebraic rank
\DeclareMathOperator{\Range}{\mathrm{Range}}	% range

\DeclareMathOperator{\Length}{\mathrm{Length}}		% length
\DeclareMathOperator{\Area}{\mathrm{Area}}			% area
\DeclareMathOperator{\Volume}{\mathrm{Vol}}			% volume
\DeclareMathOperator{\Perimeter}{\mathrm{Perim}}		% perimeter
\DeclareMathOperator{\Trees}{\mathrm{Trees}}

\DeclareMathOperator{\Volumizer}{\mathcal V}
\DeclareMathOperator{\Perimetrizer}{\mathcal P}

\newcommand{\inverse}[1]{{#1}^{-1}}		% inverse

\DeclareMathOperator{\Link}{\mathrm{Link}}		% linking number
\newcommand{\Mod}[1]{\ (\mathrm{mod}\ #1)}		% better modulus

\newcommand{\Reduced}{\widetilde{H}}						% reduced homology
\newcommand{\ReducedPersistent}{\widetilde{\mathit{PH}}}	% reduced persistent homology

\newcommand{\Normalizer}{\mathcal{Z}}		% normalization constant

\newcommand{\dualforest}[1]{{#1}^\ast}		% dual forest
\newcommand{\dual}[1]{{#1}^\bullet}			% dual subcomplex

\DeclareMathOperator{\Alexander}{\mathcal{A}}					% Alexander duality
\DeclareRobustCommand{\rAlexander}{\reflectbox{$\mathcal A$}}	% backwards cal A
\DeclareMathOperator{\Dualize}{\mathcal{D}}						% dualization map

\DeclareMathOperator{\Int}{\mathrm{Int}}	% topological interior
\DeclareMathOperator{\Bd}{\mathrm{Bd}}		% topological boundary
\DeclareMathOperator{\Cl}{\mathrm{Cl}}		% topological closure

\let\into\hookrightarrow		% injective function
\let\onto\twoheadrightarrow		% surjective function
\let\lto\longrightarrow			% longer arrow (it's prettier)
\let\lmapsto\longmapsto         % longer mapsto arrow

\DeclareMathOperator{\FILM}{\bm{\mathsf{FILM}}}		% FILM
\DeclareMathOperator{\RST}{\bm{\mathsf{RST}}}       % RST
\newcommand{\dFILM}{\FILM^*}                                 % dual FILM
\DeclareMathOperator{\CPP}{\bm{\mathsf{CPP}}}		% coupled plaquette percolation
\DeclareMathOperator{\UST}{\bm{\mathsf{UST}}}       % uniform spanning tree distribution
\DeclareMathOperator{\USF}{\bm{\mathsf{USF}}}       % uniform spanning forest distribution
\DeclareMathOperator{\PLGT}{\bm{\mathsf{PLGT}}}     % Potts lattice gauge theory
\DeclareMathOperator{\PRCM}{\bm{\mathsf{PRCM}}}     % Plaquette random cluster model
\DeclareMathOperator{\PK}{\bm{\mathsf{PK}}}         % Coupling
\DeclareMathOperator{\LERW}{\bm{\mathsf{LERW}}}		% loop-erased random walk
\DeclareMathOperator{\MST}{\bm{\mathsf{MST}}}		% minimum spanning tree
\DeclareMathOperator{\SLE}{\bm{\mathsf{SLE}}}		% SLE

\DeclareMathOperator{\Bernoulli}{\bm{\mathsf{Bernoulli}}}
\DeclareMathOperator{\RCM}{\bm{\mathsf{RCM}}}
\DeclareMathOperator{\Potts}{\bm{\mathsf{Potts}}}
\DeclareMathOperator{\ES}{\bm{\mathsf{ES}}}

\DeclareMathOperator{\H}{\mathbf{H}}        % Hamiltonian

\newcommand{\skeleton}[2]{{#1}^{(#2)}}		% skeleton

\let\free\square			% free boundary conditions
\let\wired\blacksquare		% wired boundary conditions
\let\none\boxtimes			% no boundary conditions
\let\periodic\thicksim	% periodic boundary conditions


\let\complement\overline	% complement



%%%%%%%%%%%%%%%%%%%%%%%%%%%%%%%%%%%%%%%%%%%%%%%%%%%%%%%%%%%%%%%%%%%%%%%%%%%%%%%%
%%%% SPECIAL NONMATHEMATICS TYPESETTING                                     %%%%
%%%%%%%%%%%%%%%%%%%%%%%%%%%%%%%%%%%%%%%%%%%%%%%%%%%%%%%%%%%%%%%%%%%%%%%%%%%%%%%%
\usepackage{xspace}

\let\th\undefined
\let\st\undefined
\let\nd\undefined
\let\rd\undefined

\newcommand{\th}{\textsuperscript{th}\xspace}
\newcommand{\st}{\textsuperscript{st}\xspace}
\newcommand{\nd}{\textsuperscript{nd}\xspace}
\newcommand{\rd}{\textsuperscript{rd}\xspace}

\newcommand{\qand}{\quad \text{and} \quad}
\newcommand{\qbut}{\quad \text{but} \quad}
\newcommand{\qor}{\quad \text{or} \quad}

\newcommand{\qqand}{\qquad \text{and} \qquad}
\newcommand{\qqbut}{\qquad \text{but} \qquad}
\newcommand{\qqor}{\qquad \text{or} \qquad}

\newcommand{\ATEAMS}{\textsf{ATEAMS}\xspace}
\newcommand{\ATEAMSpp}{\textsf{ATEAMS++}\xspace}


%%%%%%%%%%%%%%%%%%%%%%%%%%%%%%%%%%%%%%%%%%%%%%%%%%%%%%%%%%%%%%%%%%%%%%%%%%%%%%%%
%%%% THEMES                                                                 %%%%
%%%%%%%%%%%%%%%%%%%%%%%%%%%%%%%%%%%%%%%%%%%%%%%%%%%%%%%%%%%%%%%%%%%%%%%%%%%%%%%%

\usetheme{default}

% Set default margins.
\def\defaultbeamermargin{0.3cm}
\setbeamersize{text margin left=\defaultbeamermargin, text margin right=\defaultbeamermargin}

% Change frame title color.
\setbeamercolor{frametitle}{fg=\secondarycolor}
\setbeamerfont{frametitle}{size=\large}

% Change frame title margins.
\makeatletter
\setbeamertemplate{frametitle}[default][left,leftskip=\the\dimexpr\beamer@leftmargin-\defaultbeamermargin\relax]
\makeatother

\setbeamercolor{enumerate item}{fg=black}
\setbeamercolor{enumerate subitem}{fg=black}

% ToC colors.
\setbeamercolor{section in toc}{fg=black}

\setbeamertemplate{section in toc}[sections numbered]
\setbeamertemplate{subsection in toc}{\leavevmode\leftskip=3.2em\rlap{\hskip-2em\inserttocsectionnumber.\inserttocsubsectionnumber}\inserttocsubsection\par}

% Vacate navigation symbols.
\setbeamertemplate{navigation symbols}{}


% Change the default footline.
\setbeamercolor{section in head/foot}{%
	bg=\secondarycolor!\transparency,
	fg=black
}

\setbeamertemplate{footline}{%
	\begin{beamercolorbox}[
		wd=\paperwidth,
		ht=2.25ex,
		dp=1ex,
		leftskip=1ex,
		rightskip=1ex
	]{section in head/foot}%
		\strut%
		\ifthenelse{\equal{\insertsubsectionhead}{}}{%
			{\color{\secondarycolor}\textbf{\insertsectionhead}}
		}{%
			{\color{gray}\insertsectionhead \quad $\lto$ \quad}%
			{\color{\secondarycolor}\textbf{\insertsubsectionhead}}%
		}%
		\hfill \insertframenumber{}{\color{gray}/\inserttotalframenumber}
	\end{beamercolorbox}
}


% Inserts sub/section pages.
\AtBeginSubsection[]{
	\begin{frame}
		\begin{tikzpicture}[overlay,remember picture]
			\node[anchor=west] at (current page.west) {%
				\large{
					{\color{gray}\insertsectionnumber.\insertsubsectionnumber}%
					\quad {\color{\secondarycolor}\insertsubsectionhead}
				}
			};
		\end{tikzpicture}
	\end{frame}
}



%%%%%%%%%%%%%%%%%%%%%%%%%%%%%%%%%%%%%%%%%%%%%%%%%%%%%%%%%%%%%%%%%%%%%%%%%%%%%%%%
%%%% MISC                                                                   %%%%
%%%%%%%%%%%%%%%%%%%%%%%%%%%%%%%%%%%%%%%%%%%%%%%%%%%%%%%%%%%%%%%%%%%%%%%%%%%%%%%%

% https://tex.stackexchange.com/a/59769
% Patch to reduce space before/after proofs
\makeatletter
\renewenvironment{proof}[1][\proofname]{\par
	\vspace{-\topsep}% remove the space after the theorem
	\pushQED{\qed}%
	\normalfont
	\topsep0.5em \partopsep0pt %
	\trivlist
	\item[\hskip\labelsep
	    \itshape
	#1\@addpunct{.}]\ignorespaces
}{%
	\popQED\endtrivlist\@endpefalse
	\addvspace{0.25em plus 2pt} % some space after
}
\makeatother

% Reduce spacing around displaymath that DOESN'T get reset every time.
\makeatletter
\g@addto@macro \normalsize {%
	\setlength\abovedisplayskip{5pt plus 2pt minus 2pt}%
	\setlength\belowdisplayskip{5pt plus 2pt minus 2pt}%
}
\makeatother

